\documentclass[10pt,a4paper]{amsart}
\usepackage[margin=19mm]{geometry}
\usepackage{amsmath,amssymb,mathtools}
\usepackage[scr=rsfs,cal=euler]{mathalfa}
\usepackage{xfrac}
\usepackage[unicode,hidelinks,bookmarks=false]{hyperref}
\newcommand{\ie}{i.e.\ }
\newcommand{\cf}{cf.\ }
\newcommand{\resp}{respectively\ }
\newcommand{\Subsection}{\subsection}
\providecommand{\nobreakdash}{\nobreak\hbox{-}}
\setlength{\emergencystretch}{2em}
\multlinegap0pt
\newtheorem{theorem}{Theorem}[section]
\newtheorem{proposition}[theorem]{Proposition}
\newtheorem{definition}{Definition}[section]
\newtheorem{lemma}[theorem]{Lemma}
\newtheorem{corollary}[theorem]{Corollary}
\newtheorem{remark}{Remark}[section]
\newtheorem{example}{Example}[section]
\numberwithin{equation}{theorem}
\newcommand{\dd}[1]{\textbf{\textit{#1}}}
\title[A Note on Distance-Fall Colorings]{A Note on Distance-Fall Colorings}
\author{Wayne Goddard and Sonwabile Mafunda}
\date{Source: arXiv:2508.21232v1 (CC BY 4.0)}
\begin{document}
\begin{abstract}
We say a proper coloring  of a graph is distance-$k$ fall if every vertex is within distance $k$ 
of at least one vertex of every color. We show that 
if $G$ is a connected graph of order at least $3$ that is $3$-colorable, then
it has a distance-2 fall 3-coloring. Further, for every integer $k\ge 2$,
if $T$ is a tree of order at least $k$, then $T$ has a $k$-coloring 
such that every vertex is within distance $k-1$ of every color.
This proves an old conjecture of Beineke and Henning that every tree
of order $n$ has an independent distance-$d$-dominating set of size at most $n/(d + 1)$.
\end{abstract}
\maketitle
\noindent\emph{Source and licence.} A Note on Distance-Fall Colorings. Version arXiv:2508.21232v1 (CC BY 4.0), distributed under the Creative Commons Attribution 4.0 International licence, http://creativecommons.org/licenses/by/4.0/, as recorded in the arXiv OAI licence metadata for this exact version and shown on the abstract page https://arxiv.org/abs/2508.21232v1. Full rights notice: licenses/arxiv-2508.21232-CC-BY-4.0.txt.\par\smallskip
\noindent\textit{Editorial scope.} Counted unit: the note's first and main theorem, source label t:threeGood (source0.tex lines 119--121), printed as Theorem 2.1 in the published version and as Theorem 2.1 here, delivered together with its complete original author proof at source lines 123--146 as the single primary proof body. The proof is a run-in block, not an environment: the note contains no \textbackslash{}begin\{proof\} anywhere, and the counted proof opens with the author's own printed marker \{\bf Proof:\} at source line 123 and closes with \textbackslash{}hfill\$\textbackslash{}Box\$ at source line 146; the wrapper adds no proof marker of its own and prints exactly the author's. No mathematics is written, repaired or re-derived: the proof body is the source's own text, including its own typographic forms and line breaks. The note is an article (source0.tex line 14) with authblk affiliations, and the wrapper is the lane's pinned amsart editorial wrapper; the note's own \textbackslash{}newtheorem declarations (source lines 27--38, including the separate section-numbered definition, remark and example counters and \textbackslash{}numberwithin\{equation\}\{theorem\}) and its own macro \textbackslash{}dd (source line 38) are copied verbatim into the wrapper preamble, so every printed number of the delivered unit is produced by the note's own counters and coincides with the published version: Section 1 Introduction, Section 2 Distance-Fall Colorings and Chromatic Number, and the counted Theorem 2.1. Required context retained uncounted: the whole Introduction (source lines 85--104), which defines a fall coloring, a distance-k fall coloring, an independent set, a dominating set, a distance-k dominating set and the independent distance-k domination number and points to the 2000 paper of Dunbar et al.\ and to Beineke and Henning, and the opening of Section 2 (source lines 111--117), namely the section heading with its label, the remark that the bipartite colouring of a connected nontrivial bipartite graph is a fall colouring, and the definitions of d-good and d-bad together with the lead-in sentence that announces the counted theorem. Every definition, remark and citation of those source lines is retained. The note's abstract (source lines 62--71) is reproduced verbatim in the wrapper front matter and the note's own title and author names are reproduced in the wrapper title block. Not part of this unit and not redistributed: the note's affiliation block, keywords and MSC line (source lines 49--53 and 77--78); the rest of Section 2 (source lines 148--166, the discussion of independent isolating sets and the lead-in to the tree theorem); the tree theorem t:treeGood with its run-in proof (source lines 168--186) and the general partial-colouring theorem t:threeGoodGeneral with its run-in proof (source lines 203--220); the closing comments on graph operations (source lines 229--242); the Acknowledgment (source lines 249--250); and the note's bibliography listing, of which only the cited entries are transcribed. The delivered text cites two keys, BH and DHHJKLR, and those two entries are transcribed verbatim from the note's own in-source thebibliography (source lines 257--284) with the note's printed bibliography numbers preserved (BH prints 1, DHHJKLR prints 5), so the printed citation numbers of the delivered unit are the published ones. Reference adaptation: none was needed and none was made. The retained spans contain no cross-reference command at all: the note's five \textbackslash{}ref\{t:threeGood\} occurrences (source lines 154, 161, 188, 199 and 216) all lie outside the retained spans, so the only reference machinery inside the delivered unit is the note's own section label dfc\&cn and its theorem label t:threeGood, both of which resolve inside the unit and produce no unresolved reference. Printed slips kept literal: no printed word, symbol, hypothesis or number of the counted statement or of its proof was altered; the proof keeps the note's own \{\bf Proof:\} marker, its proof-end form (the source line 146 ends with \textbackslash{}hfill\$\textbackslash{}Box\$ followed by a line break), its spelling of ``colour/coloring'', its doubled spaces and its use of the notations N(v), N[v] and \textbackslash{}deg without definition. The note's cite package and its pgf, tikz, pgfplots, float, graphicx, tikz-cd, comment and authblk loads are not needed by the retained spans (the note has no figure, no table, no \textbackslash{}includegraphics, no \textbackslash{}input and no non-ASCII character) and are not reproduced; the note's \textbackslash{}usepackage\{cite\} only reorders multiple-key citations, and every citation of the delivered text is a single-key citation, so no printed citation changes. No mathematics is written, repaired or re-derived; all mathematics of the unit is native text with no image dependency.
\section{Introduction}
A \dd{fall coloring} is a proper coloring of the vertices of a graph such that every vertex sees every color. 
That is, for each vertex $v$ its closed neighborhood $N[v]$ contains all colors.
Not all graphs have a fall coloring. The simplest example is the $5$-cycle.
In this paper we consider proper colorings where every vertex is ``near'' to every color.
We say a coloring is \dd{distance-$k$ fall} if every vertex is within distance $k$ 
of at least one vertex of every color. For example, 
for a graph of diameter $2$, every proper coloring is distance-$2$ fall;
and every odd cycle has a $3$-coloring that is distance-$2$ fall.

The terminology ``\dd{fall coloring}'' was introduced in the 2000 paper by Dunbar et al.~\cite{DHHJKLR}. 
But the concept is older, having been studied before as partitioning a graph into independent dominating sets;
see the references of~\cite{DHHJKLR}. Recall that a set $S$ of vertices is \dd{independent} if no two elements of $S$ are joined by an edge.
A set $S$ is \dd{dominating} if every vertex is either in $S$ or adjacent to at least one vertex of $S$.
An \dd{independent dominating} set is one that is both independent and dominating.
More generally, a set $S$ is \dd{distance-$k$ dominating} if every vertex is within distance $k$ of at least
one vertex of $S$. Thus a distance-$k$ fall coloring is a partition of the vertex set into independent distance-$k$ dominating sets.
The parameter \dd{independent distance-$k$ domination}, that is, the  
minimum size of an independent distance-$k$ dominating set,
was originally studied by Beineke and Henning~\cite{BH}.
%%%%%%%%%%%%%%%%%%%%%%%%%%%%%%%%%%%%%%%%%%%%%%%%%%%%%%%%%%%%%%%%%%%%%%%%%%%%%%%%%%%%%%%%%%%%%%%%%%%%%%%%%%%%%%%%%%%%%%%%%%%%%%%%%%%%





 \section{Distance-Fall Colorings and Chromatic Number}\label{dfc&cn}

It is immediate that  if a (connected nontrivial) graph is bipartite, then the bipartite coloring is a fall coloring.
So the next question to consider is $3$-colorable graphs.
For a coloring, we say a vertex is \dd{$d$-good} if every color appears within distance $d$ of the vertex; otherwise the vertex
is \dd{$d$-bad}. The following theorem shows that if the graph is $3$-colorable then 
 there is a proper coloring with $3$ colors such that every vertex is $2$-good.

\begin{theorem} \label{t:threeGood}
If $G$ is a connected graph of order at least $3$ that is $3$-colorable, then $G$ has a distance-$2$ fall $3$-coloring.
\end{theorem}

{\bf Proof:} 
Suppose to the contrary, that $G$ does not admit a distance-$2$ fall $3$-coloring. 
Take a proper $3$-coloring of $G$ that minimizes the number of $2$-bad vertices. 
Without loss of generality, choose a $2$-bad vertex $v \in V(G)$. 
Then $N(v)$ is monochromatic.

If there exists a vertex $w_1 \in N(v)$ with $\deg(w_1)=1$, 
then by recoloring $w_1$ with the missing color in $N[v]$, 
we obtain a proper $3$-coloring of $G$ with strictly fewer $2$-bad vertices, 
contradicting the minimality of the original coloring.

Therefore, assume $\deg(w) \geq 2$ for all $w \in N(v)$. 
Since $v$ is $2$-bad, we have $N(v) \cap N(w) = \emptyset$ and $N(w)$ is monochromatic for every $w \in N(v)$. 
That is, $v$ and all vertices in its second neighborhood share the same color.

Now recolor $v$ with the missing color in $N[v]$. 
Vertex $v$ then becomes $2$-good, and every neighbor of $v$ remains $2$-good. 
Moreover, no previously $2$-good vertex becomes $2$-bad under this recoloring. 
Thus, the total number of $2$-bad vertices decreases, 
contradicting the minimality of the original coloring.

Therefore, there exists a proper $3$-coloring of $G$ with zero $2$-bad vertices.

\hfill$\Box$\\

\begin{thebibliography}{99}
\bibitem[1]{BH} Beineke, L.; Henning, M. A.; \textit{Some extremal results on independent distance domination in graphs.} Ars Combin. {\bf 37} (1994), 223--233.
\bibitem[5]{DHHJKLR} Dunbar, J. E. ; Hedetniemi, S. M.; Hedetniemi, S. T .; Jacobs, D. P.; Knisely, J.; Laskar, R. C.; Rail, D. F.; \textit{Fall Colorings of Graphs.} J. Combin. Math. Combin. Comput {\bf 33} (2000), 257--273.
\end{thebibliography}
\end{document}
